\uuid{EyOj}
\titre{Saut de la fonction de répartition}
\auteur{Maxime Nguyen}
\organisation{AMSCC}
\datecreate{2026-11-24}
\niveau{L2}
\module{Probabilité}
\chapitre{Probabilité continue}
\sousChapitre{Espérance et Fonction de répartition}
\difficulte{2}
\theme{fonction de répartition, loi mixte, probabilité d'un point}

\contenu{
\begin{question}
Une variable aléatoire réelle $X$ a pour fonction de répartition $F(t) = \mathbb{P}(X \leq t)$ :
$$F(t) = \begin{cases} 0 & \text{si } t < 0,\\ 0{,}35\,t & \text{si } 0 \leq t < 2,\\ 1 & \text{si } t \geq 2.\end{cases}$$
On note $F(2^-) = \lim\limits_{t \to 2,\, t < 2} F(t)$ la limite à gauche de $F$ en $2$. Quelles affirmations sont vraies ?
\begin{answers}
\bad{$X$ est une variable à densité.}
\bad{$X$ est une variable discrète.}
\good{$\mathbb{P}(X = 2) = 0{,}3$.}
\good{$\mathbb{P}(X \leq 1) = 0{,}35$.}
\bad{Pour toute variable aléatoire réelle $Y$, $\mathbb{P}(Y = 2) = 0$.}
\end{answers}
\begin{explanations}
A : la fonction de répartition d'une variable à densité est continue, or $F$ saute en $2$.

B : la fonction de répartition d'une variable discrète est en escalier, or $F$ croît continûment sur $[0;2[$. $X$ n'est ni discrète ni à densité : sa loi est mixte.

C : la hauteur du saut est $\mathbb{P}(X = 2) = F(2) - F(2^-) = 1 - 0{,}7 = 0{,}3$.

D : $\mathbb{P}(X \leq 1) = F(1) = 0{,}35$.

E : faux, $X$ elle-même est un contre-exemple.

Contre-sens visés : croire que $\mathbb{P}(X = a) = 0$ pour toute
variable réelle ; croire que toute variable est soit discrète, soit à densité.
\end{explanations}
\end{question}
}
